% a06_WN_b: part 2 of 4 of the walled reading of Ba3 Prop. 1 (input from sec/a06_BALtorus_d7b.tex): classes closed under the moves, the two families of the rotation part, and the statement.

\paragraph{Classes closed under the moves.} Recall the moves (P), (P$'$), (W), (X) and the classes of \cref{def:a07:class}, and the cut terms and orbits of \cref{def:a07:cut}. A \emph{decorated graph} of the walled model is a graph of the model (lines $G^v_s(T_\eta,0)$ and the torus vector propagator, vertices \cite[(1.6)--(1.15)]{Ba3} with their labels (Loc), counterterm vertices with the model's own coefficients, finite \cite[(1.7)]{Ba3} vertices with bounded coefficients) whose external legs carry external fields which may be kernels of several variables \cite[p.~419]{Ba3}, in particular the two ends of a cut line; these legs are the \emph{slots} of the graph. A finite set $\cC$ of decorated graphs with a common slot structure (the same slots, with the same kinds, differentiations, labels and fields), a common order, a common label multiset and a common multiset of finite-vertex coefficients is \emph{closed under the moves} if every move of the set $M_0$, consisting of (P) and (P$'$) at a divergent two-scalar-leg subgraph of a member, (W) at a divergent two-vector-leg subgraph, and (X) at a \emph{bridgeless} divergent $\phi\phi A$ subgraph (at the member itself included, when the member is of one of these types), gives a member. The set $M_0$ is a subset of the moves of \cref{def:a07:class} and has the properties (A1)--(A5) of that definition, which are properties of the individual moves; so \cref{lem:a07:orb} applies with $M=M_0$. Per-label full classes are closed under the moves (\cref{def:a07:class}), and so is every class $\Cl_{M}(K)$, $M\supset M_0$, of a connected component $K$ of a cut graph in \cref{lem:a07:orb}(c), with the cut ends as slots. The restriction to bridgeless subgraphs in (X) is forced by the two-point settings: the exchange at a bridged divergent $\phi\phi A$ subgraph carrying an external leg produces a one-particle-reducible graph, which is not a term of a family of \cite[(1.23)]{Ba3} graphs (\cref{thm:a06:WN}(iii)(c) shows that no partner is needed there). In a \emph{two-point setting} the members have two external scalar legs, undifferentiated, one of them, the \emph{root}, at a vertex pinned at the point $x$, the other at a vertex that is summed; the coefficient is a function of $x$, and the bounds below hold uniformly in $x$. A $\phi\phi A$ subgraph $S$ is \emph{bridgeless} if no internal line of $S$ disconnects it (a self-loop is not a bridge), and \emph{bridged} otherwise.

\paragraph{The two families of the rotation part.} (a) \emph{The cut expansion of the finite order-$4$ mass counterterm.} The order-$4$ two-point coefficient of the walled model is computed once with the $\eps$-level lines $C^v_0$, $C^\gamma$ and once with the step-$K$ lines $G^v_K$, $G^\gamma_K$; the difference, with the scalar lines decomposed as $C^v_0=G^v_K+R^v_K$ and the vector lines as $C^\gamma=G^\gamma_K+R^\gamma_K$, is the sum \eqref{eq:a07:tel} of \cref{def:a07:cut} over the cut terms with at least one cut line, the cut ends carrying the kernels $R^v_K$, $R^\gamma_K$. By \cref{lem:a07:orb}(c),(d) each orbit is a product of classes $\Cl_M(K)$ of the connected components $K$ of its cut graph, each closed under the moves, with the cut ends as slots; the component carrying the root is a two-point class with the root pinned at $x$, the others are vacuum components with slots. (b) \emph{The line-derivative family of the wall vertex.} Let $\cF^{(4)}$ be the set of two-point graphs of \cite[(1.23)]{Ba3} of coupling order $4$: one-particle-irreducible, two undifferentiated external scalar legs, the root at $x$, the inner \cite[(1.7)]{Ba3} vertices carrying the model's own step-$s$ counterterm of their order and their non-renormalising parts as finite vertices. For $\Gamma\in\cF^{(4)}$ and a scalar line $l$ of $\Gamma$ let $\Gamma_{l\to v}$ be $\Gamma$ with $l$ subdivided by a two-scalar-leg vertex $v$ at a point $z$ with a bounded coefficient $c(z)$ (the wall vertex), and
\[
(\Gamma_{l\to v})_{\mathrm{ren}}\ \defeq\ \sum_{F}\prod_{\gamma\in F}(-T_\gamma)\,\Gamma_{l\to v},
\]
the sum over the forests $F$ of divergent two-scalar-leg subgraphs of $\Gamma$ avoiding $l$ (i.e.\ not containing $l$), $T_\gamma$ the counterterm operation of (Agr). The family is $\sum_{\Gamma\in\cF^{(4)}}\sum_l(\Gamma_{l\to v})_{\mathrm{ren}}$, with no cut line.

\begin{lemma}[the line-derivative family]\label{lem:a06:WNNC}
\textup{(a)} Every connected subgraph of $\Gamma_{l\to v}$ that contains $v$ has degree $D\ge1$ or vanishes. Hence the divergent subgraphs of $\Gamma_{l\to v}$ are exactly the divergent subgraphs of $\Gamma$ avoiding $l$, and $(\Gamma_{l\to v})_{\mathrm{ren}}$ is the renormalised forest sum of $\Gamma_{l\to v}$. \textup{(b)} A forest of divergent two-scalar-leg subgraphs of $\Gamma\in\cF^{(4)}$ avoiding $l$ has at most one element. \textup{(c)} Group the terms $T_F\Gamma_{l\to v}\defeq\prod_{\gamma\in F}(-T_\gamma)\Gamma_{l\to v}$ ($\Gamma\in\cF^{(4)}$, $l$ a scalar line of $\Gamma$, $F$ as above) by label multiset and by the multiset of finite-vertex coefficients, the wall vertex included. Each group is a two-point class closed under the moves, with the root pinned at $x$.
\end{lemma}

\begin{proof}
(a) The vertex $v$ is a \cite[(1.7)]{Ba3} vertex with a bounded coefficient, so it has $D(v)=2$ and every connected subgraph $H$ containing it has $D(H)\ge1$ if $H$ has at least two slots (\cref{lem:a06d:S25}(b), \cite[Cor.~2.3]{Ba3}). A connected subgraph with exactly one slot vanishes: with a scalar slot by the symmetry $\phi\to-\phi$ of every vertex (its amputated kernel $K_a$ with one open scalar index satisfies $K_a=-K_a$), with a vector slot at zero background because its scalar index loops close inside it and $\Tr q^{\mathrm{odd}}=0$ (Step~3(D$_4$)); and a connected subgraph of $\Gamma_{l\to v}$ with no slot is $\Gamma_{l\to v}$ itself, which has two. The divergent subgraphs of $\Gamma_{l\to v}$ therefore avoid $v$, i.e.\ are subgraphs of $\Gamma$ not containing $l$, with the same vertices, lines, slots and degree (\cref{lem:a07:orb}(a)); conversely a divergent subgraph of $\Gamma$ avoiding $l$ is a subgraph of $\Gamma_{l\to v}$ with the same data.

(b) A proper divergent two-scalar-leg subgraph of $\Gamma$ has coupling order $2$ (\cite[(1.23)]{Ba3}; the order-$4$ one is $\Gamma$ itself, which contains $l$). Two such subgraphs cannot be nested, and two disjoint ones $\gamma,\gamma'$ would exhaust the vertices of $\Gamma$: each has exactly two slots, the two external legs of $\Gamma$ are among the four slots, so at most one line joins $\gamma$ to $\gamma'$; if exactly one does it is a bridge of $\Gamma$, and if none does $\Gamma$ is disconnected; both contradict $\Gamma\in\cF^{(4)}$.

(c) Fix a group. \emph{(P), (P$'$).} A divergent two-scalar-leg subgraph $\gamma$ of a member $T_F\Gamma_{l\to v}$ avoids $v$ and $l$ by (a); the move replaces the term with forest $F$ by the term with forest $F\cup\{\gamma\}$ (or $F\setminus\{\gamma\}$), a forest avoiding $l$ by (b), of the same $\Gamma$ and $l$; the counterterm vertex carries $\gamma$'s labels (Agr) and the model's own coefficient, so the labels and the finite-vertex coefficients of the group are kept. \emph{(W).} A divergent two-vector-leg subgraph $H$ of a member avoids $v$ and $l$ by (a). Replacing $H$ by a graph $H'$ of its per-label two-vector-leg class (same two vector slots, same order, same label multiset, its counterterm members included) changes nothing outside $H$: the line $l$ with the vertex $v$ on it and the two external scalar legs are untouched, the order and the labels are kept, and the result is one-particle-irreducible: since $T_F\Gamma_{l\to v}$ is one-particle-irreducible, the two attachment points of $H$ are connected outside $H$ (else $H$'s vector line to the rest would be a bridge), so an internal bridge of $H'$ separating its two slots is not a bridge of the new graph, and an internal bridge of $H'$ with a slot-free side bounds a one-slot subgraph, which vanishes by (a). Smoothing $v$ back into a line gives a graph of $\cF^{(4)}$ with the counterterm vertices of $F$ inside or outside $H$ as the case may be (a forest element inside $H$ is a counterterm member of $H$'s class), so the result is a member of the group. \emph{(X).} The exchange at a bridgeless $\phi\phi A$ subgraph $P$ moves only the attachments at $P$'s two scalar slots and no vertex: the line $l$ keeps its ends up to this re-attachment and $v$ stays on it, the external legs stay undifferentiated (both slots of $P$ carry no difference, since $P$ is divergent), and the result is one-particle-irreducible: for a line $\ell\notin P$, $P$ is connected in both graphs minus $\ell$ with the same multiset of outer attachment points, and for $\ell\in P$, $P\setminus\ell$ is connected because $P$ is bridgeless. The forest element, if any, is nested with $P$ or disjoint from it (a divergent two-scalar-leg subgraph crossing a bridgeless order-$3$ $\phi\phi A$ subgraph would meet it in a single one-vector-leg vertex whose two internal lines of $P$ would be its two scalar slots, leaving its vector line as a bridge of the forest element); a forest element sharing a vertex with $P$ but no line (the tadpole at the \cite[(1.6)]{Ba3} vertex of the double-line shape) is untouched, since \textup{(X)} moves attachments only; so $F$ is carried along. \emph{Bridged $\phi\phi A$ subgraphs.} No move of $M_0$ acts there. (The exchange at a bridged divergent $\phi\phi A$ subgraph $\{u\}\cup S_2$ of a member, with the external leg of the member at $u$, would move that leg to $S_2$, which then hangs on the bridge alone: the result is one-particle-reducible and is not a term of the family. Such a member exists: $u$, $u'$ one-vector-leg vertices \cite[(1.8)]{Ba3} joined by a vector line, $w$ a \cite[(1.6)]{Ba3} vertex with a self-loop, scalar lines $u$--$w$ and $w$--$u'$, the root at $u$ and the wall vertex on $w$--$u'$; the node $\{u,w\}$ is a dangerous bridged $\phi\phi A$ node of degree $0$.) Finally every member has its root pinned at $x$, since no move moves a vertex.
\end{proof}

\begin{theorem}[the walled form of {\cite[Prop.~1]{Ba3}} for classes closed under the moves]\label{thm:a06:WN}
In the walled model at a step $s\le K$ (potential $0\le v\le\bar v$ on the scalar lines only, vector lines unwalled, zero background, the model's own counterterms, finite \cite[(1.7)]{Ba3} vertices with bounded coefficients, torus aligned with the blocks with sides $\ge1$) let $\cC$ be a class closed under the moves, of order $\le\bar n$, with slot fields $\Phi_{\mathrm{ext}}$, $A_{\mathrm{ext}}$ and finite-vertex coefficients $c_v$. Then:
\begin{enumerate}[label=\textup{(\roman*)},nosep,leftmargin=*]
\item\textup{(regrouping)} \Cref{lem:a07:orb} (with $M=M_0$ in the two-point settings, under which the families of \cite[(1.23)]{Ba3} graphs are closed) and \cref{thm:a06d:S23} hold verbatim with walled lines; in particular the cut expansion of the finite order-$4$ mass counterterm is a finite sum over orbits, each orbit the contraction through the kernels $R^v_K$, $R^\gamma_K$ of per-component classes closed under the moves, with the invariants of \cref{lem:a07:orb}(b),(f); and the line-derivative family of the wall vertex is a finite disjoint union of two-point classes closed under the moves (\cref{lem:a06:WNNC}).
\item\textup{(the bound)} With $e_s=e(L^s\eps)$ and $\lambda_s=\lambda(L^s\eps)$,
\[
\big|E^v(\cC;\{\square(v)\},\Phi_{\mathrm{ext}},A_{\mathrm{ext}})\big|\ \le\ O(1)^v\,e_s^{\,d_v}\lambda_s^{\,d_s}\prod_{v\ \mathrm{finite}}\sup_z|c_v(z)|\;e^{-\delta^v_0d(\{\square(v)\})}\prod\norm{h\Phi_{\mathrm{ext}}}_{1,\alpha_0}\prod\norm{h'A_{\mathrm{ext}}}_{1,\alpha_0},
\]
the form \cite[(1.33)]{Ba3} ($E^v$ the class sum of \cref{lem:a07:orb}(d) with walled lines): $d_v$, $d_s$ count the couplings of the action vertices, a counterterm vertex counting as the exact counterterm of the order it replaces; the localisations $h$, $h'$ are those of (H$_A$); the slot fields are kernels of several variables, each measured by its norm \cite[(1.32)]{Ba3} in the variable at the slot; in a two-point setting the left side is the coefficient at the pinned root $x$, and the bound is uniform in $x$. The constants $O(1)^v$ and $\delta^v_0$ depend only on the order, on the parameters of Step~3(D) ($d=3$, $L=5$, $a$, $N$, $\norm q$, $\bar n$, $\alpha_0$, $m$, $\lmu$, $M$, $R$) and polynomially on $\bar v$; they do not depend on $\eps$, on $\eps_0\in(0,\eps_0^*]$ below the caps of Step~3, on $s$, $K$, the labels, the torus side, $x$, or on $v$ beyond $\bar v$.
\item\textup{(the $\phi\phi A$ node clauses)} Let $S$ be a dangerous $\phi\phi A$ node of a member $G$ of $\cC$ along an ordering, and let $c$ be its regrouped zeroth-order local coefficient in the form of \cref{thm:a06d:S48}: rescaled to $\xi\Z^3$, lines $\breve G_l=C^\xi+R_l$ with $R_l=R^v_l$ of \cref{lem:a06:WN1}(d) on the scalar lines and the bulk remainder on the vector lines, multipliers at the anchor, and the decomposition $c=c^C+c^R$ into the all-$C^\xi$ part and the terms with at least one $R$-line.
\begin{enumerate}[label=\textup{(\alph*)},nosep,leftmargin=*]
\item $S$ bridgeless, $G^{\mathrm{sw}}$ the member obtained by \textup{(X)} at $S$, a common anchor in $S$: $c(S\text{ in }G)+c(S\text{ in }G^{\mathrm{sw}})=0$ exactly, for any line kernels.
\item $S$ bridgeless with its two scalar slots at distinct vertices $u_1\ne u_2$, the pinned root at $u_1$ in $G$ and at $u_2$ in $G^{\mathrm{sw}}$, anchors at the root; $c_{S,1}(x)$, $c_{S,2}(x)$ the two coefficients: $c^C_{S,1}(x)=c^C_{S,2}(x)$, and $|c_{S,1}(x)-c_{S,2}(x)|\le C^v_7$.
\item $S$ bridged. Then $S=\{u\}\cup S_2$, where $u$ is a one-vector-leg vertex \cite[(1.8)]{Ba3} carrying one scalar slot and the vector slot of $S$, its difference acting on the bridge $l_0=(u,w)$, $w\in S_2$, and $S_2$ is a two-scalar-leg subgraph of order $2$ (the exchange bubble, the \cite[(1.6)]{Ba3} tadpole, a seagull with a vector loop, or a \cite[(1.7)]{Ba3} vertex); these are the shapes \cite[(2.21b)--(g)]{Ba3}. If $u$ is summed, then with the anchor in $S_2$, $c=0$ for every member and any line kernels. If $u$ carries the pinned root (anchor $u$ at $x$), then $c^C=0$ for every member, and the $R$-part of the pair $(S\text{ with }S_2)+(S\text{ with }T_{S_2})$, both members of $\cC$ by \textup{(P)}, is $\le C^v_{\mathrm{BR}}$ in absolute value; the $R$-part of a single member is in general not bounded uniformly in $\xi$.
\end{enumerate}
$C^v_7$ and $C^v_{\mathrm{BR}}$ are the $\phi\phi A$ constant of \cref{thm:a06d:S48} times a polynomial in $\bar v$, free of $\xi$, $j_0$, $s$, $\eps$, the labels, the torus and $x$.
\end{enumerate}
\end{theorem}

The rotation part consumes (ii) and (iii) at the step $s=K$ for the per-component classes of the cut expansion in (i), and at every step $s=a\le K$ for the line-derivative family, both with $v$ in the family $v_\theta$ and at coupling order $\le4$. The proof of (ii) is by induction on the number of vertices of the members, as the proof of \cite[(1.33)]{Ba3} is (\cref{lem:a06d:S25}); the induction runs over classes closed under the moves, the completed pieces of a cut partnered subgraph being such classes of lower order.

\begin{proof}[Proof of \textup{(i)}]
\Cref{lem:a07:orb} uses the lines only through the properties (A1)--(A5) of the moves, the multilinearity of the expansion \eqref{eq:a07:tel} and the symmetry of the decomposed kernels; $R^v_K=C^v_0-G^v_K$ is a difference of two symmetric resolvents, hence symmetric. In the two-point setting the family of \cite[(1.23)]{Ba3} graphs is closed under $M_0$: (P) contracts or expands a one-particle-irreducible two-leg subgraph (a bridge of the new graph would be one of the old), and (W), (X) at bridgeless subgraphs keep one-particle-irreducibility and the undifferentiated external legs by the arguments given in the proof of \cref{lem:a06:WNNC}(c). \Cref{thm:a06d:S23} and \cref{lem:a06d:S22} are statements about labels and the cut rule. Degrees count lines (CV). So (i) is the bulk statement, read with walled lines, and the second half is \cref{lem:a06:WNNC}(c).
\end{proof}
